% ch03_u05 —— 原书 p.101-p.111（3.7.5 末句至 3.7.17 末；本单元不含节标题）
\SourcePage{101}

我们不证明这个定理。事实上，唯一不平凡的陈述是：（99）（等价地，（34））是环同态；证明见 ???。
% [Tutor Note: 原书 p.101 此处交叉引用印作「???」（引用未解析），且与下文之间缺空格；底本如此，保留原样。]

上述定理的自然途径，基于 VOA（即顶点算子代数）这一概念，或基于它的几何版本；后者在 [BD] 中以「手性代数」之名引入。\footnote{在 2.9.4 -- 2.9.5 中，我们用到了 VOA 理论（或手性代数理论）的一些想法。} 下一小节（读者可以跳过）概述手性代数的途径。

\bdp{3.7.6} 光滑曲线 $X$ 上的手性代数，是一个（左）$\mathcal{D}_X$-模 $\mathcal{A}$，带一个态射
\begin{equation}\tag{101}
j_*j^{!}(\mathcal{A}\boxtimes\mathcal{A})\to\Delta_*\mathcal{A},
\end{equation}
其中 $\Delta:X\hookrightarrow X\times X$ 是对角线，$j:(X\times X)\setminus\Delta(X)\hookrightarrow X$。
% [Tutor Note: 原书 p.101 此处把 $j$ 的靶写成 $X$；按定义应为 $(X\times X)\setminus\Delta(X)\to X\times X$，底本如此，保留原样。]

态射（101）应满足若干公理，这里不予陈述。若（101）把 $\mathcal{A}\boxtimes\mathcal{A}$ 映为 0，则称该手性代数是\emph{交换的}。此时（101）诱导一个态射 $\Delta_*(\mathcal{A}\otimes\mathcal{A})=j_*j^{!}(\mathcal{A}\boxtimes\mathcal{A})/\mathcal{A}\boxtimes\mathcal{A}\to\Delta_*\mathcal{A}$；换言之，诱导一个态射
\begin{equation}\tag{102}
\mathcal{A}\otimes\mathcal{A}\to\mathcal{A}.
\end{equation}

此时手性代数公理只要求：带运算（102）的 $\mathcal{A}$ 是交换结合含幺代数。于是，交换手性代数与 2.6 意义下的交换结合含幺 $\mathcal{D}_X$-代数相同。另一方面，由 2.6.5，对应于 $\operatorname{Aut}O$-模 $\mathit{Vac}'$ 的 $\mathcal{D}_X$-模 $\mathit{Vac}'_X$ 具有自然的手性代数结构（见下面的注）。由嵌入 $\mathfrak{z}_{\mathfrak{g}}(O)\to\mathit{Vac}'$ 诱导的映射 $\mathfrak{z}_{\mathfrak{g}}(O)_X\to\mathit{Vac}'_X$ 是手性代数态射。给定点 $x\in X$，可定义从手性代数到结合拓扑代数的函子 $\mathcal{A}\mapsto\mathcal{A}_{((x))}$。若对某个交换 $\operatorname{Aut}O$-代数 $A$ 有 $\mathcal{A}=A_X$，则 $\mathcal{A}_{((x))}$ 是 3.7.3 中的代数 $A_{K_x}$。若 $\mathcal{A}=\mathit{Vac}'_X$，则 $\mathcal{A}_{((x))}$ 是完备扭泛包络代数 $\overline{U}'=\overline{U}'(\mathfrak{g}\otimes K)$。于是

\SourcePage{102}

由函子性，得到态射 $\mathfrak{z}_{\mathfrak{g}}(K)=\mathfrak{z}_{\mathfrak{g}}(O)_K\to\overline{U}'$。其像含于 $\mathfrak{Z}$ 中，因为 $\mathfrak{z}_{\mathfrak{g}}(O)_X$ 是手性代数 $\mathit{Vac}'_X$ 的\emph{中心}。

\emph{注.} 我们概述 $\mathit{Vac}'_X$ 上手性代数结构的定义。首先，对每个 $n$，在 $\operatorname{Sym}^nX$ 上构造一个 $\mathcal{D}$-模 $\mathit{Vac}'_{\operatorname{Sym}^nX}$（$n=1$ 时得到 $\mathit{Vac}'_X$）。$\mathit{Vac}'_{\operatorname{Sym}^nX}$ 在 $D\in\operatorname{Sym}^nX$ 处的纤维 $\mathit{Vac}'_D$ 可描述如下。把 $D$ 视为 $X$ 的 $n$ 阶闭子概形，记 $O_D$ 为 $X$ 沿 $D$ 的形式完备化上的函数环，并以 $\operatorname{Spec}K_D=(\operatorname{Spec}O_D)\setminus D$ 定义 $K_D$。与 $n=1$ 的情形一样，定义 $\mathfrak{g}\otimes K_D$ 的中心扩张 $\widetilde{\mathfrak{g}\otimes K_D}$。$\mathit{Vac}'_D$ 是对应于 Harish-Chandra 对 $(\widetilde{\mathfrak{g}\otimes K_D},G(O_D))$ 的扭真空模（见 1.2.5）。记 $\mathit{Vac}'_{X\times X}$ 为 $\mathit{Vac}'_{\operatorname{Sym}^2X}$ 到 $X\times X$ 的拉回。则
\begin{equation}\tag{103}
j^{!}\mathit{Vac}'_{X\times X}=j^{!}(\mathit{Vac}'_X\boxtimes\mathit{Vac}'_X),
\end{equation}
\begin{equation}\tag{104}
\Delta^{\dagger}\mathit{Vac}'_{X\times X}=\mathit{Vac}'_X
\end{equation}
其中 $j$ 与 $\Delta$ 的含义同（101），$\Delta^{\dagger}$ 表示朴素拉回，即 $\Delta^{\dagger}=H^1\Delta^{!}$。把（101）定义为复合
\[
j_*j^{!}\mathit{Vac}'_X\boxtimes\mathit{Vac}'_X=j_*j^{!}\mathit{Vac}'_{X\times X}\to j_*j^{!}\mathit{Vac}'_{X\times X}/\mathit{Vac}'_{X\times X}=\Delta_*\mathit{Vac}'_X
\]
最后一步等式来自（104）。

\bdp{3.7.7} \emph{定理.}
\begin{enumerate}[label=(\roman*),leftmargin=2.2em]
\item 态射（100）是拓扑同构。
\item $G(K)$ 在 $\mathfrak{Z}$ 上的伴随作用是平凡的。
\end{enumerate}
证明将在 3.7.10 中给出。它基于 Feigin--Frenkel 定理，故下面两点是本质的：$\mathfrak{g}$ 半单；$\mathfrak{g}\otimes K$ 的中心扩张对应「临界」数量积（18）。对定理 3.7.5 而言，这两点并不本质。

我们还将证明以下结论。

\bdp{3.7.8} \emph{定理.} 2.9.8 中定义的映射 $\operatorname{gr}\mathfrak{Z}\to\mathfrak{Z}^{cl}$ 诱导拓扑同构
\[
\operatorname{gr}_i\mathfrak{Z}\xrightarrow{\ \sim\ }\mathfrak{Z}^{cl}_{(i)}:=\{\mathfrak{Z}^{cl}\text{ 中 }i\text{ 次齐次多项式构成的空间}\}.
\]

\SourcePage{103}

\bdp{3.7.9} \emph{定理.} 记 $\mathcal{I}_n$ 为 $\overline{U}'$ 中由 $\mathfrak{g}\otimes t^nO$（$n\ge 0$）拓扑生成的闭左理想。则理想 $I_n:=\mathcal{I}_n\cap\mathfrak{Z}\subset\mathfrak{Z}$ 由空间 $\mathfrak{Z}^m_i$（$m<i(1-n)$）拓扑生成；这里 $\mathfrak{Z}^m_i:=\{z\in\mathfrak{Z}_i\mid L_0z=mz\}$，$\mathfrak{Z}_i$ 是 $\mathfrak{Z}$ 的标准滤链，$L_0:=-t\frac{d}{dt}\in\operatorname{Der}O$。

\bdp{3.7.10} 我们来证明上述各定理。像（100）的元素是 $G(K)$-不变的（见 2.9.6 的注）。于是 3.7.7(ii) 由 3.7.7(i) 推出。下面证明 3.7.7(i)、3.7.8 与 3.7.9。

由 2.5.2，$\operatorname{gr}\mathfrak{z}_{\mathfrak{g}}(O)=\mathfrak{z}^{cl}_{\mathfrak{g}}(O)$。按 2.4.1，$\mathfrak{z}^{cl}_{\mathfrak{g}}(O)$ 可与 $\mathfrak{g}^*\otimes\omega_O$ 上 $G(O)$-不变多项式函数构成的环等同。选取 $\mathfrak{g}^*$ 上 $G$-不变多项式代数的齐次生成元 $p_1,\dots,p_r$，并置 $d_j:=\deg p_j$。如下定义 $v_{jk}\in\mathfrak{z}^{cl}_{\mathfrak{g}}(O)$（$1\le j\le r$，$0\le k<\infty$）：
\begin{equation}\tag{105}
p_j(\eta)=\sum_{k=0}^{\infty}v_{jk}(\eta)t^k(dt)^{d_j},\qquad \eta\in\mathfrak{g}^*\otimes\omega_O.
\end{equation}

按 2.4.1，代数 $\mathfrak{z}^{cl}_{\mathfrak{g}}(O)$ 由 $v_{jk}$ 自由生成。$\operatorname{Der}O$ 在 $\mathfrak{z}^{cl}_{\mathfrak{g}}(O)$ 上的作用容易描述。特别地 $v_{jk}=(L_{-1})^kv_{j0}/k!$，$L_0v_{j0}=d_jv_{j0}$。把 $v_{j0}\in\mathfrak{z}^{cl}_{\mathfrak{g}}(O)=\operatorname{gr}\mathfrak{z}_{\mathfrak{g}}(O)$ 提升为元素 $u_j\in\mathfrak{z}_{\mathfrak{g}}(O)$，使得 $L_0u_j=d_ju_j$。于是代数 $\mathfrak{z}_{\mathfrak{g}}(O)$ 由 $u_{jk}:=(L_{-1})^ku_j/k!$（$1\le j\le r$，$0\le k<\infty$）自由生成。与 3.7.3 末尾的例子一样，我们看到 $\mathfrak{z}_{\mathfrak{g}}(O)_K=\mathbb{C}[[\dots,\tilde{u}_{j,-1},\tilde{u}_{j0},\tilde{u}_{j1},\dots]]$，且 $L_0\tilde{u}_{jk}=(d_j+k)\tilde{u}_{jk}$。

记 $\bar{u}_{jk}$ 为 $\tilde{u}_{jk}$ 在 $\mathfrak{Z}$ 中的像。由 2.9.8，$\bar{u}_{jk}\in\mathfrak{Z}_{d_j}$；且 $\bar{u}_{jk}$ 在 $\mathfrak{Z}^{cl}_{(d_j)}$ 中的像，是由
\begin{equation}\tag{106}
p_j(\eta)=\sum_k\tilde{v}_{jk}(\eta)t^k(dt)^{d_j},\qquad \eta\in\mathfrak{g}^*\otimes\omega_K
\end{equation}
定义的函数 $\tilde{v}_{jk}:\mathfrak{g}^*\otimes\omega_K\to\mathbb{C}$。
% [Tutor Note: 原书 p.103 对一般的 $k$ 仍写 $\bar{u}_{jk}\in\mathfrak{Z}_{d_j}$；按 $L_0\bar{u}_{jk}=(d_j+k)\bar{u}_{jk}$，其滤链次数应为 $d_j+k$，底本如此，保留原样。]

我们有拓扑代数的同构
\begin{equation}\tag{107}
\mathfrak{Z}^{cl}=\mathbb{C}[[\dots\tilde{v}_{j,-1},\tilde{v}_{j0},\tilde{v}_{j1},\dots]]
\end{equation}

\SourcePage{104}

因为
\begin{equation}\tag{108}
\parbox{0.76\linewidth}{\raggedright
$\mathfrak{g}^*\otimes t^{-n}\omega_O$ 上 $G(O)$-不变多项式函数构成的代数，由 $\tilde{v}_{jk}$（$k\ge -nd_j$）的限制自由生成；而当 $k<-nd_j$ 时，$\tilde{v}_{jk}$ 到 $\mathfrak{g}^*\otimes t^{-n}\omega_O$ 的限制等于 0。}
\end{equation}

（这一陈述可直接化归为 2.4.1 中考虑的 $n=0$ 情形。）定理 3.7.8 由（107）推出。

现在考虑态射 $f_n:\mathfrak{z}_{\mathfrak{g}}(O)_K\to\mathfrak{Z}/I_n$，其中 $I_n$ 已在 3.7.9 中定义。我们将证明
\begin{equation}\tag{109}
\parbox{0.74\linewidth}{\raggedright
$f_n$ 是满射，且其核是由 $u_{jk}$（$k<d_j(1-n)$）拓扑生成的理想 $J_n$。}
\end{equation}
定理 3.7.7 与 3.7.9 由（109）推出。

为证明（109），考虑复合 $\bar{f}_n:\mathfrak{z}_{\mathfrak{g}}(O)_K\to\mathfrak{Z}/I_n\hookrightarrow(\overline{U}'/\mathcal{I}_n)^{G(O)}$。给 $\overline{U}'/\mathcal{I}_n$ 配上由 $\overline{U}'$ 上标准滤链诱导的滤链。$L_0$ 在此滤链第 $i$ 项上的特征值 $\ge i(1-n)$。故 $\operatorname{Ker}\bar{f}_n\supset J_n$，其中 $J_n$ 已在（109）中定义。现在 $\operatorname{gr}(\overline{U}'/\mathcal{I}_n)^{G(O)}$ 含于 $(\operatorname{gr}\overline{U}'/\mathcal{I}_n)^{G(O)}$ 中，后者即 $\mathfrak{g}^*\otimes t^{-n}\omega_O$ 上 $G(O)$-不变多项式的代数。利用（108）容易看出：由 $\bar{f}_n$ 诱导的映射 $\mathfrak{z}_{\mathfrak{g}}(O)_K/J_n\to(\overline{U}'/\mathcal{I}_n)^{G(O)}$ 是同构。这蕴含（109）。我们还证明了
\begin{equation}\tag{110}
\parbox{0.55\linewidth}{\raggedright
映射 $\mathfrak{Z}\to(\overline{U}'/\mathcal{I}_n)^{G(O)}$ 是满射}
\end{equation}
从而
\begin{equation}\tag{111}
\mathfrak{Z}=(\overline{U}')^{G(O)}.
\end{equation}

\bdp{3.7.11} \emph{注}
\begin{enumerate}[label=(\roman*),leftmargin=2.2em]
\item 这里给出（111）的另一个证明。\footnote{它类似于下述事实的证明：$\mathfrak{g}\otimes K$ 的可积离散表示是平凡的。我们无法直接使用这一事实，因为 $\overline{U}'$ 不离散。} 设 $u\in(\overline{U}')^{G(O)}$。取 $h\in H(K)$，其中 $H\subset G$ 是固定的嘉当子群。则 $h^{-1}uh$ 关于某个 Borel 子群 $B_h\subset G$ 不变。
\SourcePage{105}
于是 $h^{-1}uh$ 是 $G$-不变的（只需对 $h^{-1}uh$ 在\emph{离散空间} $\overline{U}'/\mathcal{I}_n$ 中的像证明即可，其中 $\mathcal{I}_n$ 已在 3.7.9 中定义）。因此，对任意 $h\in H(K)$，$u$ 关于 $h\mathfrak{g}h^{-1}\subset\mathfrak{g}\otimes K$ 不变。李代数 $\mathfrak{g}\otimes K$ 由 $\mathfrak{g}\otimes O$ 与 $h\mathfrak{g}h^{-1}$（$h\in H(K)$）生成。故 $u\in\mathfrak{Z}$。
\item 事实上，对任意开子代数 $\mathfrak{a}\subset\mathfrak{g}\otimes K$，都有
\begin{equation}\tag{112}
\mathfrak{Z}=(\overline{U}')^{\mathfrak{a}}.
\end{equation}
确实，可按如下方式修改上述证明。把 $u$ 写成 $u_\chi$ 的（无限）和，其中 $\chi\in\mathfrak{h}^*:=(\operatorname{Lie}H)^*$，且对 $a\in\mathfrak{h}$ 有 $[a,u_\chi]=\chi(a)u_\chi$。然后取 $h\in H(K)$，使得 $h$ 在 $H(K)/H(O)=\{\text{余权格}\}$ 中的像，关于包含 $\mathfrak{h}$ 的 Borel 子代数 $\mathfrak{b}\subset\mathfrak{g}$ 是「非常占优」的，从而 $h^{-1}\mathfrak{a}h\supset[\mathfrak{b},\mathfrak{b}]$。我们看到：除非 $\chi$ 占优，$u_\chi=0$；且 $h^{-1}u_0h$ 是 $\mathfrak{g}$-不变的。用 $h^{-1}$ 替换 $h$，可见 $u=u_0$，等等。
\item 这里给出 3.7.7(ii) 的另一个证明。考虑 $\overline{U}'$ 的典范滤链 $\overline{U}'_k$。由（109）可知，空间 $\overline{U}'_k\cap\mathfrak{Z}$（$k\in\mathbb{N}$）之并在 $\mathfrak{Z}$ 中稠密。故只需证明：对每个 $k$，$G(K)$ 在 $\overline{U}'_k\cap\mathfrak{Z}$ 上的作用是平凡的。$G(K)$ 在 $\mathfrak{Z}^{cl}$ 上的作用是平凡的（见（107）、（106））。故 $G(K)$ 在 $\operatorname{gr}\mathfrak{Z}$ 上的作用是平凡的。$\mathfrak{g}\otimes K$ 在 $\widetilde{\mathfrak{g}\otimes K}$ 上对应于由（19）定义的 $G(K)$ 作用的那个作用，就是伴随作用；而 $\mathfrak{g}\otimes K$ 在 $\mathfrak{Z}$ 上的伴随作用是平凡的。故 $G(K)$ 在 $\mathfrak{Z}$ 上的作用经由 $\pi_0(G(K))$ 分解。群 $\pi_0(G(K))$ 有限（见 4.5.4），证毕。
\end{enumerate}

\bdp{3.7.12} 我们将从 [FF92] 推出定理 3.6.7。记 $A_{{}^{L}\mathfrak{g}}(O)$ 为 $\operatorname{Op}_{{}^{L}\mathfrak{g}}(O)$ 的坐标环（即 $\operatorname{Spec}O$ 上 ${}^{L}\mathfrak{g}$-oper 构成的概形）。设 $\varphi_O:A_{{}^{L}\mathfrak{g}}(O)\xrightarrow{\ \sim\ }\mathfrak{z}_{\mathfrak{g}}(O)$ 是满足 3.2.2 中条件的同构。它诱导 $\operatorname{Aut}K$-等变同构 $\varphi_K:A_{{}^{L}\mathfrak{g}}(K)\xrightarrow{\ \sim\ }\mathfrak{z}_{\mathfrak{g}}(K)$，其中 $A_{{}^{L}\mathfrak{g}}(K)$ 是 3.7.3 中对应于 $A=A_{{}^{L}\mathfrak{g}}(O)$ 的代数 $A_K$。

\SourcePage{106}

回忆：$A_K$ 是 $\operatorname{Spec}A_{Y'}$ 的水平截面构成的归纳概形的坐标环，$Y':=\operatorname{Spec}K$。若 $A=A_{{}^{L}\mathfrak{g}}(O)$，则 $\operatorname{Spec}A_{Y'}$ 是 $Y'$ 上 ${}^{L}\mathfrak{g}$-oper 的 jets 概形，其水平截面是 $Y'$ 上的 ${}^{L}\mathfrak{g}$-oper（参看 3.3.3）。于是 $A_{{}^{L}\mathfrak{g}}(K)$ 是 $\operatorname{Op}_{{}^{L}\mathfrak{g}}(K)$ 的坐标环；后者定义为 $\operatorname{Spec}K$ 上 ${}^{L}\mathfrak{g}$-oper 构成的归纳概形。它是关于 \emph{Gelfand--Dikii 括号}的泊松代数（我们在 3.7.14 中回忆其定义）。Gelfand--Dikii 括号依赖于 ${}^{L}\mathfrak{g}$ 上非退化不变双线性形式的选取。我们把它取为 $\mathfrak{g}$ 上形式（18）的\emph{对偶}：即，它在 $\mathfrak{h}^*={}^{L}\mathfrak{h}\subset{}^{L}\mathfrak{g}$ 上的限制，对偶于（18）在 $\mathfrak{h}$ 上的限制。

由 3.7.5 与 3.7.7，有典范同构 $\mathfrak{z}_{\mathfrak{g}}(K)\xrightarrow{\ \sim\ }\mathfrak{Z}$，故可把 $\varphi_K$ 视为 $\operatorname{Aut}K$-等变同构
\begin{equation}\tag{113}
A_{{}^{L}\mathfrak{g}}(K)\xrightarrow{\ \sim\ }\mathfrak{Z}.
\end{equation}
$\mathfrak{Z}$ 是关于 Hayashi 括号的泊松代数（见 3.6.2）。

\bdp{3.7.13} \emph{定理.} [FF92] 存在同构
\begin{equation}\tag{114}
\varphi_O:A_{{}^{L}\mathfrak{g}}(O)\xrightarrow{\ \sim\ }\mathfrak{z}_{\mathfrak{g}}(O)
\end{equation}
满足 3.2.2 的条件，且相应的同构（113）与泊松结构相容。

我们将证明（见 3.7.16）：具备定理中所述性质的同构（114）满足 3.6.7 的条件。故它是唯一的（见 3.6.7 的注）。

\emph{注.} 如 3.6.12 所解释，可对 $\mathfrak{g}$ 上任一不变双线性形式 $B$ 配上 $\mathfrak{Z}$ 上的一个泊松括号（3.6.2 的括号对应形式（18））。若 $B$ 非退化，则可考虑 ${}^{L}\mathfrak{g}$ 上的对偶形式，以及 $A_{{}^{L}\mathfrak{g}}(K)$ 上相应的 Gelfand--Dikii 括号。对应于（114）的同构（113）与这些泊松括号相容。

\SourcePage{107}

\bdp{3.7.14} 我们回忆 [DS85] 中 Gelfand--Dikii 括号的定义。这是 $\operatorname{Op}_{\mathfrak{g}}(K)$ 上的泊松括号（即其坐标环 $A_{\mathfrak{g}}(K)$ 上的泊松括号）。它依赖于 $\mathfrak{g}$ 上非退化不变双线性形式 $(\,,\,)$ 的选取。

记 $\widetilde{\mathfrak{g}\otimes K}$ 为 $\mathfrak{g}\otimes K$ 对应于 $(\,,\,)$ 的 Kac--Moody 中心扩张。作为向量空间，$\widetilde{\mathfrak{g}\otimes K}=(\mathfrak{g}\otimes K)\oplus\mathbb{C}$；其换位子由 2-余循环 $\operatorname{Res}(du,v)$（$u,v\in\mathfrak{g}\otimes K$）定义。拓扑对偶空间 $(\widetilde{\mathfrak{g}\otimes K})^*$ 是归纳概形。$(\widetilde{\mathfrak{g}\otimes K})^*$ 上正则函数的代数，是关于 Kirillov 括号\footnote{如 [We83] 所解释，「Kirillov 括号」由 Sophus Lie 发明，后来又被包括 A.A. Kirillov 在内的若干人重新发现。}的泊松代数（即唯一的连续泊松括号，使得从 $\widetilde{\mathfrak{g}\otimes K}$ 到 $(\widetilde{\mathfrak{g}\otimes K})^*$ 上正则函数代数的自然映射是李代数态射）。故 $(\widetilde{\mathfrak{g}\otimes K})^*$ 是泊松「流形」。记 $(\widetilde{\mathfrak{g}\otimes K})^*_1$ 为连续线性泛函 $l:\widetilde{\mathfrak{g}\otimes K}\to\mathbb{C}$ 构成的空间，其中 $l$ 在中心 $\mathbb{C}\subset\widetilde{\mathfrak{g}\otimes K}$ 上的限制是恒等。$(\widetilde{\mathfrak{g}\otimes K})^*_1$ 是 $(\widetilde{\mathfrak{g}\otimes K})^*$ 的泊松子流形。

我们把 $(\widetilde{\mathfrak{g}\otimes K})^*_1$ 与 $\operatorname{Conn}$ 等同，后者定义为 $\operatorname{Spec}K$ 上平凡 $G$-丛的联络构成的归纳概形：对联络 $\nabla=d+\eta$（$\eta\in\mathfrak{g}\otimes\omega_K$），配以 $l\in(\widetilde{\mathfrak{g}\otimes K})^*_1$，使得对任意 $u\in\mathfrak{g}\otimes K\subset\widetilde{\mathfrak{g}\otimes K}$ 有 $l(u)=\operatorname{Res}(u,\eta)$。容易验证：$\mathfrak{g}\otimes K$ 在 $\operatorname{Conn}$ 上的规范作用，对应于 $\mathfrak{g}\otimes K$ 在 $(\widetilde{\mathfrak{g}\otimes K})^*_1$ 上的余伴随作用；而 $G(K)$ 在 $(\widetilde{\mathfrak{g}\otimes K})^*$ 上的余伴随作用\footnote{它是对偶于由（19）定义的 $G(K)$ 在 $\widetilde{\mathfrak{g}\otimes K}$ 上的伴随作用（当然，（19）中的 $c$ 应换成我们 $\mathfrak{g}$ 上的双线性形式）。}，定义为：其在 $(\widetilde{\mathfrak{g}\otimes K})^*_1$ 上的限制对应于 $G(K)$ 在 $\operatorname{Conn}$ 上的规范作用。$G(K)$ 在泊松「流形」$(\widetilde{\mathfrak{g}\otimes K})^*_1$ 上的作用按字面意义不是哈密顿的，即：不能定义矩映射 $(\widetilde{\mathfrak{g}\otimes K})^*_1\to(\mathfrak{g}\otimes K)^*$。然而可以定义矩映射 $(\widetilde{\mathfrak{g}\otimes K})^*_1\to(\widetilde{\mathfrak{g}\otimes K})^*$：它就是恒等映射。

\SourcePage{108}

要点在于：$\operatorname{Op}_{\mathfrak{g}}(K)$ 可由 $\operatorname{Conn}=(\widetilde{\mathfrak{g}\otimes K})^*_1$ 经哈密顿约化得到（这种对 $\operatorname{Op}_{\mathfrak{g}}(K)$ 的解释，自动定义 $A_{\mathfrak{g}}(K)$ 上的泊松括号）。固定 Borel 子群 $B\subset G_{ad}$。设 $N$ 是其幂幺根，$\mathfrak{n}:=\operatorname{Lie}N$。由于 Kac--Moody 余循环在 $\mathfrak{n}\otimes K$ 上的限制平凡，有显然的分裂 $\mathfrak{n}\otimes K\to\widetilde{\mathfrak{g}\otimes K}$。它是 $B(K)$-等变的，且此性质唯一刻画它。$N(K)$ 在 $\operatorname{Conn}$ 上的作用是哈密顿作用：矩映射 $\mu:\operatorname{Conn}=(\widetilde{\mathfrak{g}\otimes K})^*_1\to(\mathfrak{n}\otimes K)^*$ 由上述分裂诱导。设 $\operatorname{Char}^{\star}\subset(\mathfrak{n}\otimes K)^*$ 是\emph{非退化特征标}的集合，即满足下述条件的李代数态射 $l:\mathfrak{n}\otimes K\to\mathbb{C}$ 之集：对每个单根 $\alpha$，$l$ 在 $\mathfrak{g}^\alpha\otimes K$ 上的限制非零。对每个 $l\in\operatorname{Char}^{\star}$，$N(K)$ 在 $\mu^{-1}(l)$ 上的作用自由，且商 $N(K)\backslash\mu^{-1}(l)$ 可典范地等同于 $\operatorname{Op}_{\mathfrak{g}}(K)$：确实，$\mu^{-1}(l)$ 由满足下述条件的联络 $\nabla=d+\eta\in\operatorname{Conn}$ 构成：$\eta=\sum_{\alpha\in\Gamma}J_\alpha+q$，其中 $q\in\mathfrak{b}\otimes\omega_K$，$\Gamma$ 是单根集，$J_\alpha=J_\alpha(l)$ 是 $\mathfrak{g}^{-\alpha}\otimes\omega_K$ 中固定的非零元。于是，$\operatorname{Op}_{\mathfrak{g}}(K)$ 由 $\operatorname{Conn}$ 对 $N(K)$ 的作用作关于 $l$ 的哈密顿约化得到，由此得到 $\operatorname{Op}_{\mathfrak{g}}(K)$ 上的泊松括号。它称为\emph{Gelfand--Dikii 括号}。它不依赖于 $l$。确实，对 $l,l'\in\operatorname{Char}^{\star}$，考虑同构
\begin{equation}\tag{115}
N(K)\backslash\mu^{-1}(l)\xrightarrow{\ \sim\ }N(K)\backslash\mu^{-1}(l')
\end{equation}
它来自（115）两端与 $\operatorname{Op}_{\mathfrak{g}}(K)$ 的等同。$H(K)$ 在 $\operatorname{Conn}=(\widetilde{\mathfrak{g}\otimes K})^*_1$ 上的（余）伴随作用保持相关结构（即 $\operatorname{Conn}$ 上的泊松括号、$N(K)$ 在 $\operatorname{Conn}$ 上的作用，以及矩映射 $\mu:\operatorname{Conn}\to(\mathfrak{n}\otimes K)^*$）。存在唯一的 $h\in H(K)$ 把 $l$ 变为 $l'$，且（115）由该 $h$ 的作用诱导。故（115）是泊松映射。

\emph{注}
\begin{enumerate}[label=(\roman*),leftmargin=2.2em]
\item 若 $\mathfrak{g}$ 上的双线性形式 $(\,,\,)$ 乘以 $c\in\mathbb{C}^*$，则 $\operatorname{Op}_{\mathfrak{g}}(K)$ 上的泊松括号乘以 $c^{-1}$。
\SourcePage{109}
\item 上面定义的 Gelfand--Dikii 括号是「第二 Gelfand--Dikii 括号」。第一个括号的定义，以及它与 Gelfand--Dikii 原始工作（[GD76]、[GD78]）之关系的说明，可在 [DS85] 中找到（见同上书第 2.3、3.6、3.7、6.5、8 节）。
\end{enumerate}

\bdp{3.7.15} 设 $\mathfrak{F}\in\operatorname{Op}_{\mathfrak{g}}(K)$，即 $\mathfrak{F}=(\mathfrak{F}_B,\nabla)$，其中 $\mathfrak{F}_B$ 是 $\operatorname{Spec}K$ 上的 $B$-丛，$\nabla$ 是相应 $G$-丛上满足 3.1.3 条件的联络（这里 $G$ 是对应于 $\mathfrak{g}$ 的伴随群，$B\subset G$ 是 Borel 子群）。我们将描述切空间 $T_{\mathfrak{F}}\operatorname{Op}_{\mathfrak{g}}(K)$ 与余切空间 $T^*_{\mathfrak{F}}\operatorname{Op}_{\mathfrak{g}}(K)$。然后写出 $\{\varphi,\psi\}(\mathfrak{F})$（$\varphi,\psi\in A_{\mathfrak{g}}(K)$）的显式公式。

\emph{注.} 当然 $\mathfrak{F}_B$ 总是平凡的，故可把 $\mathfrak{F}$ 视为平凡 $G$-丛中的联络 $\nabla$（即 $\nabla=d+q$，$q\in\mathfrak{g}\otimes\omega_K$），模去关于 $B$ 的规范变换。

为描述 $T_{\mathfrak{F}}\operatorname{Op}_{\mathfrak{g}}(K)$，须研究 $\mathfrak{F}=(\mathfrak{F}_B,\nabla)$ 的无穷小形变。由于 $\mathfrak{F}_B$ 不能形变，所有这些形变都来自 $\nabla$ 的无穷小形变，其形如 $\nabla(\varepsilon)=\nabla+\varepsilon q$，$q\in H^0(\operatorname{Spec}K,\mathfrak{g}^{-1}_{\mathfrak{F}}\otimes\omega_K)$（$\mathfrak{g}^{-1}$ 的定义见 3.1.1；$\mathfrak{g}^{-1}_{\mathfrak{F}}:=\mathfrak{g}^{-1}_{\mathfrak{F}_B}$ 是 $\mathfrak{g}^{-1}$ 的 $\mathfrak{F}_B$-扭）。计入 $\mathfrak{F}_B$ 的无穷小自同构，得到：
\begin{equation}\tag{116}
T_{\mathfrak{F}}\operatorname{Op}_{\mathfrak{g}}(K)=H^0(\operatorname{Spec}K,\operatorname{Coker}(\nabla:\mathfrak{b}_{\mathfrak{F}}\to\mathfrak{g}^{-1}_{\mathfrak{F}}\otimes\omega_K)).
\end{equation}
这里有切空间更方便的描述：
\begin{equation}\tag{117}
T_{\mathfrak{F}}\operatorname{Op}_{\mathfrak{g}}(K)=\operatorname{Coker}(\nabla:\mathfrak{n}^K_{\mathfrak{F}}\to\mathfrak{b}^K_{\mathfrak{F}}\otimes\omega_K)
\end{equation}
其中 $\mathfrak{n}^K_{\mathfrak{F}}:=H^0(\operatorname{Spec}K,\mathfrak{n}_{\mathfrak{F}})$，$\mathfrak{b}^K_{\mathfrak{F}}:=H^0(\operatorname{Spec}K,\mathfrak{b}_{\mathfrak{F}})$（从（117）右端到（116）右端的自然映射是同构）。利用 $\mathfrak{g}$ 上的不变数量积 $(\,,\,)$，把 $\mathfrak{b}^*$、$\mathfrak{n}^*$ 与 $\mathfrak{g}/\mathfrak{n}$、$\mathfrak{g}/\mathfrak{b}$ 等同，得到余切空间的如下描述：
\begin{equation}\tag{118}
T^*_{\mathfrak{F}}\operatorname{Op}_{\mathfrak{g}}(K)=\{u\in\mathfrak{g}^K_{\mathfrak{F}}\mid\nabla(u)\in\mathfrak{b}^K_{\mathfrak{F}}\otimes\omega_K\}/\mathfrak{n}^K_{\mathfrak{F}}.
\end{equation}

\SourcePage{110}

Gelfand--Dikii 括号的显式公式如下：
\begin{equation}\tag{119}
\{\varphi,\psi\}(\mathfrak{F})=\operatorname{Res}(\nabla(d_{\mathfrak{F}}\varphi),d_{\mathfrak{F}}\psi),\qquad \varphi,\psi\in A_{\mathfrak{g}}(K).
\end{equation}
此公式中，微分 $d_{\mathfrak{F}}\varphi$ 与 $d_{\mathfrak{F}}\psi$ 视为（118）右端的元素。

\bdp{3.7.16} \emph{定理.}\footnote{受 [Phys] 启发。}
\begin{enumerate}[label=(\roman*),leftmargin=2.2em]
\item 置 $I:=\operatorname{Ker}(A_{\mathfrak{g}}(K)\to A_{\mathfrak{g}}(O))$。则 $\{I,I\}\subset I$，故 $I/I^2$ 是 $A_{\mathfrak{g}}(O)$ 上的李代数胚。
\item 存在李代数胚的 $\operatorname{Aut}O$-等变拓扑同构
\begin{equation}\tag{120}
I/I^2\xrightarrow{\ \sim\ }\mathfrak{a}_{\mathfrak{g}}
\end{equation}
（$\mathfrak{a}_{\mathfrak{g}}$ 的定义见 3.5.11、3.5.15）。
\end{enumerate}
（此定理中，$I^2$ 表示由 $ab$（$a\in I$，$b\in I$）生成的子空间的\emph{闭包}。）

定理 3.6.7 由 3.7.13 与 3.7.16 推出。

\emph{注.} 同构（120）是唯一的（见 3.5.13 或 3.5.14）。

\bdp{3.7.17} 我们来证明定理 3.7.16。沿用 3.7.15 的记号。取 $\mathfrak{F}\in\operatorname{Op}_{\mathfrak{g}}(O)$。与（117）类似，$T_{\mathfrak{F}}\operatorname{Op}_{\mathfrak{g}}(O)$ 有如下描述：
\begin{equation}\tag{121}
T_{\mathfrak{F}}\operatorname{Op}_{\mathfrak{g}}(O)=\operatorname{Coker}(\nabla:\mathfrak{n}^O_{\mathfrak{F}}\to\mathfrak{b}^O_{\mathfrak{F}}\otimes\omega_O)
\end{equation}
其中 $\mathfrak{n}^O_{\mathfrak{F}}:=H^0(\operatorname{Spec}O,\mathfrak{n}_{\mathfrak{F}})$。$I/I^2$ 在 $\mathfrak{F}$ 处的纤维是余法空间 $T^{\perp}_{\mathfrak{F}}\operatorname{Op}_{\mathfrak{g}}(O)\subset T^*_{\mathfrak{F}}\operatorname{Op}_{\mathfrak{g}}(K)$。按（121），它用（118）的语言有如下描述：
\begin{equation}\tag{122}
T^{\perp}_{\mathfrak{F}}\operatorname{Op}_{\mathfrak{g}}(O)=\{u\in\mathfrak{g}^O_{\mathfrak{F}}\mid\nabla(u)\in\mathfrak{b}^O_{\mathfrak{F}}\otimes\omega_O\}/\mathfrak{n}^O_{\mathfrak{F}}.
\end{equation}

现在 $\{I,I\}\subset I$ 是清楚的：若 $\varphi,\psi\in I$，$\mathfrak{F}\in\operatorname{Op}_{\mathfrak{g}}(O)$，则 $d_{\mathfrak{F}}\varphi$ 与 $d_{\mathfrak{F}}\psi$ 属于（122）的右端，从而（119）的右端等于 0。

\SourcePage{111}

映射
\begin{equation}\tag{123}
I/I^2\to\operatorname{Der}A_{\mathfrak{g}}(O),
\end{equation}
是 $I/I^2$ 上代数胚结构的一部分，由 $\varphi\mapsto\partial_\varphi$ 定义：$\partial_\varphi(\psi):=\{\varphi,\psi\}$，$\varphi\in I$，$\psi\in A_{\mathfrak{g}}(K)/I=A_{\mathfrak{g}}(O)$。于是按（119），由（123）诱导的映射
\begin{equation}\tag{124}
T^{\perp}_{\mathfrak{F}}\operatorname{Op}_{\mathfrak{g}}(O)\to T_{\mathfrak{F}}\operatorname{Op}_{\mathfrak{g}}(O)
\end{equation}
是算子
\begin{equation}\tag{125}
\nabla:\{u\in\mathfrak{g}^O_{\mathfrak{F}}\mid\nabla(u)\in\mathfrak{b}^O_{\mathfrak{F}}\otimes\omega_O\}/\mathfrak{n}^O_{\mathfrak{F}}\to(\mathfrak{b}^O_{\mathfrak{F}}\otimes\omega_O)/\nabla(\mathfrak{n}^O_{\mathfrak{F}}).
\end{equation}

$I/I^2$ 上的代数胚结构，在映射（124）的核 $\mathfrak{a}_{\mathfrak{F}}$ 上诱导李代数结构。另一方面，$\mathfrak{a}_{\mathfrak{F}}$ 是（125）的核，即 $\mathfrak{a}_{\mathfrak{F}}=\{u\in\mathfrak{g}^O_{\mathfrak{F}}\mid\nabla(u)=0\}/\{u\in\mathfrak{n}^O_{\mathfrak{F}}\mid\nabla(u)=0\}$。由于 $\{u\in\mathfrak{n}^O_{\mathfrak{F}}\mid\nabla(u)=0\}=0$，有
\begin{equation}\tag{126}
\mathfrak{a}_{\mathfrak{F}}=\{u\in\mathfrak{g}^O_{\mathfrak{F}}\mid\nabla(u)=0\}.
\end{equation}
（126）的右端是 $\mathfrak{g}^O_{\mathfrak{F}}$ 的李子代数。

\emph{引理.} 由 $I/I^2$ 上代数胚结构给出的 $\mathfrak{a}_{\mathfrak{F}}$ 上的李代数结构，与由（126）诱导者一致。

\emph{证明.} 只需证明：若 $\varphi_1,\varphi_2\in A_{\mathfrak{g}}(K)$ 且 $d_{\mathfrak{F}}\varphi_i\in\mathfrak{a}_{\mathfrak{F}}$，则
\begin{equation}\tag{127}
d_{\mathfrak{F}}\{\varphi_1,\varphi_2\}=[d_{\mathfrak{F}}\varphi_1,d_{\mathfrak{F}}\varphi_2]
\end{equation}
（在（127）的右端，$d_{\mathfrak{F}}\varphi_i$ 经（126）视为 $\mathfrak{g}^O_{\mathfrak{F}}$ 的元素）。考虑 $\mathfrak{F}$ 的形变 $\mathfrak{F}(\varepsilon)$，$\varepsilon^2=0$。把 $\mathfrak{F}$ 写成 $(\mathfrak{F}_B,\nabla)$。不失一般性，可设 $\mathfrak{F}(\varepsilon)=(\mathfrak{F}_B,\nabla+\varepsilon q)$，$q\in\mathfrak{b}^K_{\mathfrak{F}}\otimes\omega_K$。把 $d_{\mathfrak{F}(\varepsilon)}\varphi_i$ 写成 $d_{\mathfrak{F}}\varphi_i+\varepsilon\mu_i$。则
\[
\begin{gathered}
\{\varphi_1,\varphi_2\}(\mathfrak{F}(\varepsilon))=\operatorname{Res}((\nabla+\varepsilon\operatorname{ad}q)(d_{\mathfrak{F}}\varphi_1+\varepsilon\mu_1),d_{\mathfrak{F}}\varphi_2+\varepsilon\mu_2)=\\
\varepsilon\operatorname{Res}([q,d_{\mathfrak{F}}\varphi_1],d_{\mathfrak{F}}\varphi_2)=\varepsilon\operatorname{Res}(q,[d_{\mathfrak{F}}\varphi_1,d_{\mathfrak{F}}\varphi_2])
\end{gathered}
\]
（我们用到了 $\nabla(d_{\mathfrak{F}}\varphi_i)=0$）。（127）的等式随之成立。\hfill$\Box$
